% set at the 10 pt body size (the article class would set the abstract environment in 9 pt)
\begin{center}\textbf{Abstract}\end{center}
\noindent Variant-effect predictors are a common source of evidence for missense variants of uncertain significance
(VUS), and the Bayesian form of the ACMG/AMP clinical guidelines expresses evidence strength as a likelihood ratio.
Gene-specific calibrations estimate this ratio from few labelled variants, as point estimates or resampling-based bounds
that are not designed to remain valid once a threshold has been chosen from the same data. We present
\textsc{CertiVar}, which turns a predictor score into loss-of-function (LoF) calls, functional calls or abstentions with
two class-conditional conformal thresholds learned from saturation genome editing or deep mutational scanning data,
keeping one variant per residue because substitutions at a residue are dependent. We prove that exact bounds on the
likelihood ratio of each call type with respect to the assay outcome remain valid after a threshold is selected from a
grid, by a Bonferroni correction or by fixed-sequence testing, and that per-class coverage after selection is at least
one minus the largest grid level. Because every
guarantee conditions on the assay class, none requires the share of LoF variants among the variants being called,
provided their within-class score distribution matches that of the assayed variants; a mixture-consistency check tests
a necessary condition for this on unlabelled targets. On 14,477 assayed missense variants in BRCA1, BRCA2, BAP1, RAD51C and MSH2, scored by
AlphaMissense, REVEL and ESM-2, the lower of the two per-class coverages was 0.896--0.908 at a 0.90 target in all 15
gene--predictor pairs, with 35\% abstention on average, and stayed at 0.89 or above when the calibration LoF share was
varied from 5\% to 50\%; marginal conformal prediction calibrated at 5\% covered LoF variants at only 0.11 in its worst
pair. Against known population likelihood ratios, the certified bounds all held together in 99.7\% of 1,500 resampled data sets
(guaranteed: at least 90\%), against 57.6\% for percentile-bootstrap bounds. At the 5\% level, the mixture check rejected 4.7\% of
random splits and every tested set of unassayed VUS outside the assayed residues. The certificates
concern the assay, not pathogenicity; certified evidence was mostly supporting-level, required same-gene assay data, and
held for 88.8\% of bounds in held-out protein regions, below the nominal 95\%. Code:
\url{https://github.com/aaryavvatts-lab/certivar}

\medskip\noindent\textbf{Keywords:} variant effect prediction; conformal prediction; likelihood ratio; ACMG/AMP
evidence calibration; multiplexed assays of variant effect; variants of uncertain significance; selective
classification; multiple testing.
